\documentclass[../../../include/open-logic-section]{subfiles}

\begin{document}

\olfileid{sth}{replacement}{limofsize}
\olsection{大小限制}

或许，为替换公理提供“内在”辩护的最常见尝试源自以下观念：

\begin{enumerate}
	\item[] \limofsize. 任何一些对象，只要其数量不是太多，就能构成一个集合。
\end{enumerate}

该原则将直接为替换公理提供辩护。毕竟，由替换所形成的任何集合，其大小都不可能超过据以形成它的集合。确切地说：假设通过替换形成集合 $\funimage{\tau}{A} = \Setabs{\tau(x)}{x \in A}$；那么 $\cardle{\funimage{\tau}{A}}{A}$；因此，如果 $A$ 的元素没有多到无法构成集合，那么它们的像也不会多到无法构成 $\funimage{\tau}{A}$。

援引 \limofsize{} 来为替换公理辩护，其显而易见的困难在于，我们迄今\emph{并未}确立任何类似 \limofsize{} 的原则。此外，当我们在 \crefrange{sth:story::chap}{sth:z::chap} 中讲述集合的累积迭代概念时，也从未有任何内容\emph{暗示}过 \limofsize{}。这也正是为何 Boolos（布洛斯）曾写道：“或许可以得出结论，集合论‘背后’至少存在两种思想” \citeyearpar[p.~19]{Boolos1989}。一方面，围绕集合之累积迭代概念的那些思想旨在为 ~$\Z$ 提供辩护。另一方面，\limofsize{} 则旨在为替换公理提供辩护。

但问题不仅仅在于我们迄今为止对 \limofsize{} 保持\emph{沉默}。真正的问题在于，\limofsize{}（如刚才所表述的）似乎与集合的累积迭代概念格格不入。毕竟，它完全没有提及集合是在\emph{阶段}中形成的这一观念。

考虑到这两种“思想”（即集合的累积迭代概念和 \limofsize{}）的历史，这其实并不太令人惊讶。这两种“思想”最终归结为两种截然不同的阻止集合论悖论的方案。集合的累积迭代概念阻止 Russell（罗素）悖论的方式大致如下：\emph{我们本就不该指望罗素集存在，因为它不会在任何阶段被“形成”}。相比之下，\limofsize{} 则旨在大致通过如下说法来排除罗素集：\emph{我们本就不该指望罗素集存在，因为它会太大}。

既然如此，不妨坦率地说：作为对悖论的一种回应，\limofsize{} 还需要\emph{远为充分的}辩护。例如，考虑这个版本的罗素悖论：\emph{没有哪只哈巴狗会恰好去嗅那些不嗅自己的哈巴狗}（参见\olref[sth][story][rus]{sec}）。如果你问“为什么没有这样一只哈巴狗？”，那么“这样一只哈巴狗势必得嗅太多的哈巴狗”并不是一个好回答。那么，为什么说罗素集必定“太大”以至于无法存在，就能成为关于其不存在的好的直观解释呢？

简而言之，如果你对 \limofsize{} 的“直观”动机感到有些困惑，那是完全可以理解的。

\end{document}